(1) 临界等温线上一阶、二阶体积导数同时为零。由 \(\partial p/\partial v=-RT/(v-b)^2+2a/v^3=0\)，\(\partial^2p/\partial v^2=2RT/(v-b)^3-6a/v^4=0\)，两式相除得到 \(v_c=3b\)。回代第一式得 \(RT_c=8a/(27b)\)，再代入状态方程得 \(p_c=a/(27b^2)\)。\(a\) 的量纲为压强乘摩尔体积平方，\(b\) 是摩尔体积，因此温度、压强与体积的量纲均正确。

(2) 令 \(p_r=p/p_c,T_r=T/T_c,v_r=v/v_c\)。将 \(v_c=3b,\ RT_c=8a/(27b)\) 代入状态方程，并用 \(p_c=a/(27b^2)\) 归一化，整理得 \(p_r=8T_r/(3v_r-1)-3/v_r^2\)。取 \(T_r=v_r=1\)，立即得到 \(p_r=8/2-3=1\)，因此三者临界值全部归一到一。这里的约化方程消除了具体物质参数，只是在原始范德瓦耳斯模型的适用范围内满足对应态。

(3) 在 \(T_r=1+t,v_r=1+w\) 且 \(|t|,|w|\ll1\) 下，使用 \(1/(1+3w/2)=1-3w/2+9w^2/4-27w^3/8+\cdots\) 与 \(1/(1+w)^2=1-2w+3w^2-4w^3+\cdots\)，得到 \(p_r=1+4t-6tw+9tw^2-\frac32w^3+O(tw^3,w^4)\)。一阶 \(w\) 与二阶 \(w\) 在 \(t=0\) 时恰好相消，因此临界等温线的主导非线性是 \(-3w^3/2\)，而不是线性回复斜率。

(4) 对有正的等温压缩率的均匀单相态，机械稳定要求 \((\partial p/\partial v)_T<0\)。旋节分解边界由导数等于零决定，故 \(RT/(v-b)^2=2a/v^3\)，即 \(T_{\rm sp}(v)=2a(v-b)^2/(Rv^3)\)。在 \(v=3b\) 此线达到临界温度。若局部斜率为正，则均匀状态对密度微扰机械不稳定；在它与真实共存线之间仍可能有亚稳态。旋节分解线不等于由两相化学势相等或 Maxwell 等面积规则决定的共存线，不能将两者混同。
